\begin{figure*}[t]
\centering
\newcommand{\timg}[1]{\includegraphics[width=\tcw]{figures/images/teaser_columns/#1.jpg}}%
\newcommand{\tcol}[2]{\begin{minipage}[t]{\tcw}\centering\timg{#1}\\[-5pt]{\tiny #2}\end{minipage}}%
\newlength{\tcw}\setlength{\tcw}{0.175\textwidth}%
%
% --- Row 1 ---
\begin{minipage}[t]{0.555\textwidth}
\centering{\scriptsize\textbf{Spatial Controls}}\\[-8pt]
\rule{\linewidth}{0.3pt}\\[1pt]
\tcol{canny}{Canny}\hfill\tcol{pose}{Pose}\hfill\tcol{depth}{Depth}
\end{minipage}%
\hspace{0.02\textwidth}%
\begin{minipage}[t]{0.36\textwidth}
\centering{\scriptsize\textbf{Camera Control}}\\[-8pt]
\rule{\linewidth}{0.3pt}\\[1pt]
\tcol{from_image}{From Image}\hfill\tcol{from_video}{From Video}
\end{minipage}%
%
\vspace{8pt}
%
% --- Row 2 ---
\begin{minipage}[t]{\tcw}
\centering{\scriptsize\textbf{Motion}}\\[-8pt]
\rule{\linewidth}{0.3pt}\\[1pt]
\tcol{sparse_tracking}{Sparse Tracking}
\end{minipage}%
\hspace{0.02\textwidth}%
\begin{minipage}[t]{0.36\textwidth}
\centering{\scriptsize\textbf{Editing}}\\[-8pt]
\rule{\linewidth}{0.3pt}\\[1pt]
\tcol{cut_on_action}{Cut-on-Action}\hfill\tcol{inpainting}{Inpainting}
\end{minipage}%
\hspace{0.02\textwidth}%
\begin{minipage}[t]{\tcw}
\centering{\scriptsize\textbf{Audio-Visual}}\\[-8pt]
\rule{\linewidth}{0.3pt}\\[1pt]
\tcol{who_is_talking}{Who Is Talking}
\end{minipage}%
\hspace{0.02\textwidth}%
\begin{minipage}[t]{\tcw}
\centering{\scriptsize\textbf{Audio Only}}\\[-8pt]
\rule{\linewidth}{0.3pt}\\[1pt]
\tcol{intensity_to_waveform}{Intensity$\to$Waveform}
\end{minipage}%
%
\caption{AVControl trains each control modality as a lightweight LoRA. Each column shows control input (top) and generated output (bottom), covering spatial controls, camera trajectory, motion, editing, and audio-visual generation.}
\label{fig:teaser}
\end{figure*}


Controlling the generation process of video and audio models is essential for practical creative applications.
However, the space of possible controls is vast: different modalities carry different types of information, and the same input (such as a mask) can have entirely different meanings depending on the context.
Rather than attempting to build a single monolithic system that handles all control types, we propose AVControl, a flexible and easily extendable framework that can be rapidly adapted to new modalities, whether standard controls like depth and pose or specialized ones such as rendering Blender previews for real-time game engines (Figure~\ref{fig:teaser}).
We build on LTX-2~\cite{hacohen2026ltx2}, a joint audio-visual DiT that natively generates synchronized video and audio, making it a natural backbone for multimodal control.

The range of controls extends well beyond spatially-aligned ControlNet-style~\cite{zhang2023controlnet} modalities like depth, pose, and canny edges.
We may wish to control camera motion from a single image, or re-render an existing video at a new trajectory while preserving scene dynamics.
When audio is also considered, the space grows further: adapting acoustics to a text-described environment, synchronizing video with a reference audio track, and more.

Our approach draws inspiration from In-Context LoRA (IC-LoRA)~\cite{huang2024iclora}, where a LoRA is trained on an image model to generate composite images, such as paired images side-by-side, with a learned relationship between the panels. At inference time, one half serves as the conditioning input while the other is generated via inpainting.
However, for structural controls such as depth, this approach fails to faithfully follow the conditioning signal (Figure~\ref{fig:canvas_vs_concat}). We hypothesize that the large spatial distance between semantically corresponding positions in the concatenated layout weakens their interaction in the attention layers.
We therefore adopted an approach inspired by Flux Kontext~\cite{batifol2025fluxkontext}, providing the reference on a parallel ``canvas,'' i.e., where additional tokens in the attention layers are processed alongside the generation target.
The challenge with this parallel layout is that the model must distinguish reference tokens from generation tokens. Flux Kontext~\cite{batifol2025fluxkontext} addresses this by introducing a new Rotary Position Embedding (RoPE)~\cite{su2021rope} dimension, which requires learning entirely new positional relationships from extensive compute and large-scale curated paired data. For video, the data cost is even more prohibitive, as temporally aligned multi-view video pairs are substantially harder to curate at scale than image pairs.
Our formulation avoids both costs. LTX-2~\cite{hacohen2026ltx2} assigns a unique per-token timestep, so the model inherently distinguishes clean reference tokens from noised generation tokens (see Section~\ref{sec:method}). The only trainable component is a minimal LoRA adapter~\cite{hu2022lora} on the frozen joint audio-visual backbone. Unlike methods that require new architectural components, this minimal formulation enables faithful video structural control where direct extensions of prior methods fail.
Moreover, because reference and target interact through self-attention, the reference influence can be continuously modulated at inference time globally or locally -- a capability unavailable to channel-concatenation methods.

Because each control modality is a lightweight, independently trained LoRA, this design enables easy extension to new controls without retraining existing ones. Unlike monolithic methods such as VACE~\cite{jiang2025vace}, which train all controls jointly, adding a new control, whether a neural renderer for Blender meshes (Section~\ref{sec:modalities}) or a speech-to-ambient audio transformation (Section~\ref{sec:av}), requires only a small dataset and a short training run. The total training budget across all 13 trained modalities is ${\sim}$55K steps, less than one third of VACE's 200K-step training run.

To accelerate inference, we further propose a small-to-large control grid that reduces the reference canvas resolution for sparse controls such as camera parameters.

\vspace{0.5em}
\noindent To summarize, our contributions are:
\begin{itemize}
    \item A compute- and data-efficient framework for training per-modality control LoRAs on a parallel canvas, enabling faithful video structural control and fine-grained inference-time strength modulation.
    \item A diverse set of independently trained control modalities, from spatially-aligned controls and camera trajectory to audio-visual applications, demonstrating the framework's flexibility.
    \item A small-to-large control grid training strategy that reduces the reference canvas resolution for spatially sparse modalities, lowering latency without sacrificing control fidelity.
\end{itemize}